\begin{thebibliography}{99}
\bibitem{inf:dlmf}
F. W. J. Olver, A. B. Olde Daalhuis, D. W. Lozier, B. I. Schneider,
R. F. Boisvert, C. W. Clark, B. R. Miller, B. V. Saunders, H. S. Cohl,
and M. A. McClain (eds.),
\emph{NIST Digital Library of Mathematical Functions},
\url{https://dlmf.nist.gov/}, Chapter 10,
\S\S10.2, 10.6, 10.16, 10.17, and 10.21, in particular
Eqs.~10.2.2, 10.16.1, and 10.21.17.
\bibitem{inf:watson}
G. N. Watson, \emph{A Treatise on the Theory of Bessel Functions},
second edition, Cambridge University Press, 1944.
\bibitem{inf:BakerWustholz}
A. Baker and G. W\"ustholz,
\emph{Logarithmic Forms and Diophantine Geometry},
New Mathematical Monographs 9, Cambridge University Press, 2008,
\S1.2.
\url{https://doi.org/10.1017/CBO9780511542862}.
\bibitem{inf:gasper}
G. Gasper, Linearization of the product of Jacobi polynomials. II,
\emph{Canadian Journal of Mathematics} \textbf{22} (1970), 582--593,
Theorem 1. \url{https://doi.org/10.4153/CJM-1970-065-4}.
\bibitem{inf:siegel}
C. L. Siegel, \emph{\"Uber einige Anwendungen diophantischer Approximationen}
(1929), reprinted in U. Zannier (ed.), \emph{On Some Applications of
Diophantine Approximations}, Edizioni della Normale, 2014, pp.~81--138.
\url{https://doi.org/10.1007/978-88-7642-520-2_2}.
\bibitem{inf:lindemann}
F. Lindemann, Ueber die Zahl $\pi$,
\emph{Mathematische Annalen} \textbf{20} (1882), 213--225.
\url{https://doi.org/10.1007/BF01446522}.
\bibitem{inf:kato}
T. Kato, \emph{Perturbation Theory for Linear Operators},
second edition, Classics in Mathematics, Springer, 1995,
Chapters II and IV--VII.
\url{https://doi.org/10.1007/978-3-642-66282-9}.
\bibitem{inf:GT}
D. Gilbarg and N. S. Trudinger,
\emph{Elliptic Partial Differential Equations of Second Order},
Classics in Mathematics, Springer, 2001, Chapters 6--9.
\url{https://doi.org/10.1007/978-3-642-61798-0}.
\bibitem{inf:MorreyNirenberg}
C. B. Morrey, Jr., and L. Nirenberg,
On the analyticity of the solutions of linear elliptic systems of
partial differential equations,
\emph{Communications on Pure and Applied Mathematics}
\textbf{10} (1957), 271--290.
\url{https://doi.org/10.1002/cpa.3160100204}.
\bibitem{inf:MorreyBoundary}
C. B. Morrey, Jr., On the analyticity of the solutions of analytic
non-linear elliptic systems of partial differential equations.
Part II. Analyticity at the boundary,
\emph{American Journal of Mathematics} \textbf{80} (1958), 219--237.
\url{https://doi.org/10.2307/2372831}.
\bibitem{inf:sacksteder}
R. Sacksteder, On hypersurfaces with no negative sectional curvatures,
\emph{American Journal of Mathematics} \textbf{82} (1960), 609--630.
\url{https://doi.org/10.2307/2372973}.
\bibitem{inf:zygmund}
A. Zygmund, \emph{Trigonometric Series}, second edition, Volumes I and II,
Cambridge University Press, 1959, Vol.~II, Ch.~X, Theorem~3.13.
\bibitem{inf:banach}
S. Banach, Sur les op\'erations dans les ensembles abstraits et leur
application aux \'equations int\'egrales,
\emph{Fundamenta Mathematicae} \textbf{3} (1922), 133--181.
\url{https://doi.org/10.4064/fm-3-1-133-181}.
\bibitem{inf:BST73}
L. Brown, B. M. Schreiber, and B. A. Taylor,
Spectral synthesis and the Pompeiu problem,
\emph{Annales de l'Institut Fourier (Grenoble)} \textbf{23} (1973),
no.~3, 125--154. \url{https://doi.org/10.5802/aif.474}.
\bibitem{inf:Wil76}
S. A. Williams, A partial solution of the Pompeiu problem,
\emph{Mathematische Annalen} \textbf{223} (1976), no.~2, 183--190.
\url{https://doi.org/10.1007/BF01360881}.
\bibitem{inf:Zal92}
L. Zalcman, A bibliographic survey of the Pompeiu problem,
in \emph{Approximation by Solutions of Partial Differential Equations}
(B. Fuglede et al., eds.), NATO ASI Series C, vol.~365,
Kluwer, Dordrecht, 1992, 185--194.
\url{https://doi.org/10.1007/978-94-011-2436-2_17}.
\bibitem{inf:CLdDP26}
G. Cao-Labora and J. de Dios Pont,
Counterexamples to Schiffer's conjecture, arXiv:2608.05114 (2026).
\url{https://arxiv.org/abs/2608.05114}.
\bibitem{inf:CS26}
M. J. Colbrook and G. Stepaniants,
A computer-assisted counterexample to the planar Pompeiu and Schiffer
conjectures, arXiv:2608.01579 (2026).
\url{https://arxiv.org/abs/2608.01579}.
\bibitem{inf:Guo26}
J. Guo,
Convex counterexamples to the Schiffer and Pompeiu conjectures in
dimensions two to eighteen, arXiv:2609.35419 (2026).
\url{https://arxiv.org/abs/2609.35419}.
\bibitem{inf:doCarmoLima}
M. P. do Carmo and E. Lima,
Isometric immersions with semi-definite second quadratic forms,
\emph{Archiv der Mathematik} \textbf{20} (1969), 173--175.
\url{https://doi.org/10.1007/BF01899010}.
Reprinted in K. Tenenblat (ed.), \emph{Manfredo P. do Carmo---Selected Papers},
Springer, Berlin, Heidelberg, 2012, 43--45.
\url{https://doi.org/10.1007/978-3-642-25588-5_4}.
\end{thebibliography}
